%Data institusi

%\input{Dozentdaten}

%\renewcommand{\fachbereich}{Bidang studi}

%\renewcommand{\dozent}{Prof. Dr. . }

%Data ujian

\renewcommand{\klausurgebiet}{ }

\renewcommand{\klausurtyp}{ }

\renewcommand{\klausurdatum}{ . 20}

\klausurvorspann {\fachbereich} {\klausurdatum} {\dozent} {\klausurgebiet} {\klausurtyp}

%Data untuk tabel poin berikut

\renewcommand{\aeins}{  3  }

\renewcommand{\azwei}{  3   }

\renewcommand{\adrei}{  7  }

\renewcommand{\avier}{  4  }

\renewcommand{\afuenf}{  0  }

\renewcommand{\asechs}{  0  }

\renewcommand{\asieben}{  0  }

\renewcommand{\aacht}{  5  }

\renewcommand{\aneun}{  15  }

\renewcommand{\azehn}{  5  }

\renewcommand{\aelf}{  0  }

\renewcommand{\azwoelf}{  3  }

\renewcommand{\adreizehn}{  3   }

\renewcommand{\avierzehn}{  0  }

\renewcommand{\afuenfzehn}{  10  }

\renewcommand{\asechzehn}{  4  }

\renewcommand{\asiebzehn}{  62  }

\renewcommand{\aachtzehn}{    }

\renewcommand{\aneunzehn}{    }

\renewcommand{\azwanzig}{    }

\renewcommand{\aeinundzwanzig}{    }

\renewcommand{\azweiundzwanzig}{    }

\renewcommand{\adreiundzwanzig}{    }

\renewcommand{\avierundzwanzig}{    }

\renewcommand{\afuenfundzwanzig}{    }

\renewcommand{\asechsundzwanzig}{    }

\punktetabellesechzehn

\klausurnote

\newpage

\setcounter{section}{K}

__NOEDITSECTION__

\inputaufgabeklausurloesung
{3}

{

Definisikan istilah-istilah berikut
\zusatzklammer {yang dicetak miring} {} {}.
\aufzaehlungsechs{Suatu kurva
\stichwort {berparametrisasi panjang busur} {}
\maabbdisp {\gamma} {[a,b]} { \R^n
} {.}

}{Untuk suatu titik

\mavergleichskette
{\vergleichskette
{  P
}
{ \in }{ M
}
{  }{}
{    }{}
{   }{}
}
{}{}{}
pada suatu manifold diferensiabel $M$, yang dimaksud dengan dua kurva diferensiabel yang
\stichwort {ekuivalen secara tangensial} {} ialah
\maabbdisp {\gamma_1, \gamma_2} { I } { M
} {}
\zusatzklammer {di sini

\mavergleichskettek
{\vergleichskettek
{  0
}
{ \in }{  I
}
{  }{
}
{    }{
}
{   }{
}
}
{}{}{}
merupakan interval riil terbuka dan

\mavergleichskettek
{\vergleichskettek
{   \gamma_1(0)
}
{ = }{ \gamma_2(0)
}
{ = }{ P
}
{    }{
}
{   }{
}
}
{}{}{}} {} {.}

}{Suatu
\stichwort {penampang kontinu} {}
dari suatu pemetaan kontinu
\maabbdisp {p} {Y} {X
} {}
antara ruang-ruang topologis
\mathkor {} {X} {dan} {Y} {.}

}{\stichwort {Bentuk volume kanonik} {} pada suatu manifold Riemann berorientasi $M$.

}{\stichwort {Turunan eksterior} {}
dari suatu
\definitionsverweis {bentuk diferensial}{}{}
berderajat $k$ yang
\definitionsverweis {terdiferensialkan secara kontinu}{}{}

\mathl{\omega \in
{ \mathcal E }^{ k } (  U   )}{} pada suatu
\definitionsverweis {himpunan terbuka}{}{}

\mathl{U \subseteq \R^n}{.}

}{\stichwort {Kelengkungan seksional} {} yang berkaitan dengan vektor-vektor bebas linear

\mavergleichskette
{\vergleichskette
{  v,w
}
{ \in }{T_PM
}
{  }{
}
{    }{
}
{   }{
}
}
{}{}{}
pada suatu
\definitionsverweis {manifold Riemann}{}{}
$M$ di suatu titik

\mavergleichskette
{\vergleichskette
{ P
}
{ \in }{ M
}
{  }{
}
{    }{
}
{   }{
}
}
{}{}{.}
}

}
{

\aufzaehlungsechs{Suatu
\definitionsverweis {kurva diferensiabel}{}{}
\maabbdisp {\gamma} {[a,b]} {\R^n
} {,}
disebut
berparametrisasi panjang busur
jika

\mavergleichskette
{\vergleichskette
{ \gamma_1'(t)^2 + \cdots + \gamma_n'(t)^2
}
{ = }{ 1
}
{  }{
}
{    }{
}
{   }{
}
}
{}{}{}
berlaku untuk setiap $t$.
}{Kedua kurva
\mathkor {} {\gamma_1} {dan} {\gamma_2} {}
disebut ekuivalen secara tangensial di $P$ jika terdapat suatu lingkungan terbuka
\mathl{P \in U}{} dan suatu
\definitionsverweis {peta}{}{}
\maabbdisp {\alpha} { U } { V
} {}
dengan
\mathl{V \subseteq \R^n}{} sedemikian sehingga

\mavergleichskettedisp
{\vergleichskette
{  \left(  \alpha \circ \left(  \gamma_1 \vert_{\gamma_1^{-1} (U)}  \right)   \right)'(0)
}
{ =} { \left(  \alpha \circ \left(  \gamma_2 \vert_{\gamma_2^{-1} (U)}  \right)  \right)'(0)
}
{ } {
}
{ } {
}
{ } {
}
}
{}{}{} berlaku.
}{Yang dimaksud dengan
penampang kontinu
dari $p$ adalah suatu pemetaan kontinu
\maabb {s} {X} {Y
} {}
yang memenuhi

\mavergleichskettedisp
{\vergleichskette
{  p \circ s
}
{ =} {
\operatorname{Id}_{ X }
}
{ } {
}
{ } {
}
{ } {
}
}
{}{}{.}
}{Untuk
\mathl{P \in M}{,} misalkan
\mathl{\omega_P}{} adalah
\definitionsverweis {bentuk alternan}{}{}
pada
\mathl{T_PM}{}
\zusatzklammer {atau elemen yang bersesuaian dari \mathlk{\bigwedge^n T^*_PM}{}} {} {,}
yang memberikan nilai $1$ kepada setiap
\definitionsverweis {basis ortonormal}{}{}
yang merepresentasikan
\definitionsverweis {orientasi}{}{.}
Maka
$n$-\definitionsverweis {bentuk diferensial
}{}{}
\maabbeledisp {} { M } { \bigwedge^n T^*M
} {P } { \omega_P
} {,}
disebut bentuk volume kanonik pada $M$.
}{Bentuk $\omega$ memiliki representasi pada $U$

\mathdisp {\omega = \sum_{I \subseteq \{1 , \ldots , n \},\,  { \# \left(   I  \right) } =k } f_I dx_I} {  }

dengan
\definitionsverweis {fungsi-fungsi yang terdiferensialkan secara kontinu}{}{}
\maabbdisp {f_I} { U } {\R
} {.}
Maka turunan eksteriornya adalah bentuk diferensial berderajat $(k+1)$

\mavergleichskettedisp
{\vergleichskette
{ d \omega
}
{ =} { \sum_{I \subseteq \{1 , \ldots , n \},\,  { \# \left(   I  \right) } = k } df_I \wedge dx_I
}
{ =} { \sum_{I \subseteq \{1 , \ldots , n \},\,  { \# \left(   I  \right) } = k } \left(  \sum_{j = 1}^n { \frac{   \partial  f_I  }{  \partial  x_j  } } dx_j  \right)  \wedge dx_I
}
{ } {
}
{ } {
}
}
{}{}{.}
}{Bilangan

\mavergleichskettedisp
{\vergleichskette
{ K(v,w)
}
{ =} {  { \frac{    \left\langle R(v,w)w , v  \right\rangle  }{   \left\langle v , v  \right\rangle  \left\langle w , w  \right\rangle  -\left\langle v , w  \right\rangle^2   } }
}
{ } {
}
{ } {
}
{ } {
}
}
{}{}{,}
dengan $R$ menyatakan
\definitionsverweis {operator kelengkungan}{}{}
pada bundel tangen, disebut
\definitionswort {kelengkungan seksional}{}
yang berkaitan dengan $v,w$ di $P$.
}

 }

__NOEDITSECTION__

\inputaufgabeklausurloesung
{3}

{

Rumuskan teorema-teorema berikut.
\aufzaehlungdrei{\stichwort {Teorema tentang orientasi pada suatu hipermuka diferensiabel} {}

\mavergleichskette
{\vergleichskette
{  Y
}
{ \subseteq }{  \R^n
}
{  }{
}
{    }{
}
{   }{
}
}
{}{}{.}}{Rumus braket Lie medan-medan vektor pada suatu himpunan terbuka

\mavergleichskette
{\vergleichskette
{  U
}
{ \subseteq }{  \R^n
}
{  }{
}
{    }{
}
{   }{
}
}
{}{}{.}}{Teorema tentang
\stichwort {retraksi ke batas} {}
pada manifold.}

}
{

\aufzaehlungdrei{\faktsituation {Misalkan

\mavergleichskette
{\vergleichskette
{ W
}
{ \subseteq }{ \R^n
}
{  }{
}
{    }{
}
{   }{
}
}
{}{}{}
\definitionsverweis {terbuka}{}{,}
\maabb {h} { W } { \R
} {}
adalah suatu
\definitionsverweis {fungsi yang terdiferensialkan secara kontinu}{}{}
dan

\mavergleichskette
{\vergleichskette
{ Y
}
{ = }{ h^{-1}(c)
}
{  }{
}
{    }{
}
{   }{
}
}
{}{}{}
adalah
\definitionsverweis {serat}{}{}
di atas

\mavergleichskette
{\vergleichskette
{ c
}
{ \in }{ \R
}
{  }{
}
{    }{
}
{   }{
}
}
{}{}{,}
dengan $h$
\definitionsverweis {reguler}{}{}
di setiap titik $Y$. Misalkan $Y$
\definitionsverweis {terhubung}{}{.}}
\faktfolgerung {Maka terdapat tepat dua
\definitionsverweis {orientasi}{}{}
pada $Y$.}
\faktzusatz {}
\faktzusatz {}}{\faktsituation {Misalkan

\mavergleichskette
{\vergleichskette
{ U
}
{ \subseteq }{ \R^n
}
{  }{
}
{    }{
}
{   }{
}
}
{}{}{}
\definitionsverweis {terbuka}{}{}
dan misalkan $f_1 , \ldots , f_n, g_1 , \ldots , g_n$ adalah
\definitionsverweis {fungsi-fungsi yang dua kali terdiferensialkan secara kontinu}{}{}
pada $U$ dengan operator-operator diferensial terkait

\mavergleichskette
{\vergleichskette
{ D
}
{ = }{ \sum_{i = 1}^n f_i \partial_i
}
{  }{
}
{    }{
}
{   }{
}
}
{}{}{}
dan

\mavergleichskette
{\vergleichskette
{ E
}
{ = }{ \sum_{j = 1}^n g_j \partial_j
}
{  }{
}
{    }{
}
{   }{
}
}
{}{}{}}
\faktfolgerung {Maka berlaku

\mavergleichskettedisp
{\vergleichskette
{ D \circ E -E \circ D
}
{ =} { \sum_{j = 1}^n \left(  \sum_{i = 1}^n \left(  f_i  \partial_i g_j - g_i \partial_i f_j  \right)  \right)  \partial_j
}
{ } {
}
{ } {
}
{ } {
}
}
{}{}{.}
Secara khusus, ekspresi ini sendiri bersesuaian dengan suatu operator diferensial berorde $1$, sehingga bersesuaian dengan suatu medan vektor.}
\faktzusatz {}
\faktzusatz {}}{Misalkan $M$ adalah suatu
\definitionsverweis {manifold diferensiabel berbatas}{}{}
\definitionsverweis {kompak}{}{}
\definitionsverweis {berorientasi}{}{}

\mathl{\partial M}{} dengan
\definitionsverweis {basis topologi terhitung}{}{.}
Maka tidak terdapat
\definitionsverweis {pemetaan yang terdiferensialkan secara kontinu}{}{}
\maabbdisp {\varphi} { M } { \partial M
} {,}
yang
\definitionsverweis {pembatasannya}{}{}
pada $\partial M$ adalah identitas.}

 }

__NOEDITSECTION__

\inputaufgabeklausurloesung
{7}

{

Buktikan
\stichwort {teorema tentang penyelesaian persamaan diferensial biasa pada suatu hipermuka diferensiabel} {}

\mavergleichskette
{\vergleichskette
{ Y
}
{ \subseteq }{ \R^n
}
{  }{
}
{    }{
}
{   }{
}
}
{}{}{.}

}
{

Kita bekerja dengan struktur Euklides pada $\R^n$ dan dengan
\definitionsverweis {gradien}{}{}
$\operatorname{Grad} \,  h$ dari $h$. Menurut asumsi, pada $Y$ gradien ini tidak pernah sama dengan $0$; karena kekontinuan, sifat ini berlaku pula pada suatu lingkungan terbuka dari $Y$. Dengan memperkecil $U$ jika perlu, kita dapat mengasumsikan bahwa gradien tidak memiliki titik nol di seluruh $U$. Pada $U$, tinjau medan vektor baru $G$ yang didefinisikan oleh

\mavergleichskettedisp
{\vergleichskette
{ G(P)
}
{ =} { F(P)-  { \frac{   \left\langle  F(P) ,
\operatorname{Grad} \,  h  (  P )    \right\rangle   }{  \Vert {
\operatorname{Grad} \,  h  (  P ) } \Vert^2   } }
\operatorname{Grad} \,  h  (  P )
}
{ } {
}
{ } {
}
{ } {
}
}
{}{}{}
dengan rumus tersebut. Untuk

\mavergleichskette
{\vergleichskette
{ P
}
{ \in }{ Y
}
{  }{
}
{    }{
}
{   }{
}
}
{}{}{}
berlaku

\mavergleichskette
{\vergleichskette
{ G(P)
}
{ = }{ F(P)
}
{  }{
}
{    }{
}
{   }{
}
}
{}{}{,}
karena pada $Y$ gradien dari $h$ tegak lurus terhadap medan vektor. Selain itu, $G$
\zusatzklammer {berbeda dengan $F$} {} {}
memiliki sifat bahwa untuk setiap titik

\mavergleichskette
{\vergleichskette
{ P
}
{ \in }{ U
}
{  }{
}
{    }{
}
{   }{
}
}
{}{}{}
vektor $G(P)$ tegak lurus terhadap gradien; memang,

\mavergleichskettealignhandlinks
{\vergleichskettealignhandlinks
{ \left\langle  G(P) ,
\operatorname{Grad} \,  h  (  P )   \right\rangle
}
{ =} { \left\langle  F(P)- { \frac{   \left\langle  F(P) ,
\operatorname{Grad} \,  h  (  P )    \right\rangle  }{  \Vert {
\operatorname{Grad} \,  h  (  P ) } \Vert^2  } }
\operatorname{Grad} \,  h  (  P )  ,
\operatorname{Grad} \,  h  (  P )   \right\rangle
}
{ =} { \left\langle  F(P) ,
\operatorname{Grad} \,  h  (  P )   \right\rangle - \left\langle  F(P) ,
\operatorname{Grad} \,  h  (  P )   \right\rangle
}
{ =} { 0
}
{ } {
}
}
{}
{}{.}
Sekarang misalkan
\maabbdisp {v} { I} { \R^n
} {}
adalah, menurut
[[Picard Lindelöf/Lokal Lipschitz/Lokale Existenz und Eindeutigkeit/Fakt|Satz 56.2 (Analysis (Osnabrück 2021-2023))]]
, solusi tunggal dari masalah nilai awal untuk $G$ dengan syarat awal

\mavergleichskette
{\vergleichskette
{ v(t_0)
}
{ = }{ P
}
{ \in }{ Y
}
{    }{
}
{   }{
}
}
{}{}{.}
Maka

\mavergleichskettealign
{\vergleichskettealign
{ \left(Dh \circ v \right)_{t}
}
{ =} { \left(Dh  \right)_{v(t)} \circ  \left(D v \right)_{t}
}
{ =} { \left(Dh  \right)_{v(t)}  \left(  v'(t)  \right)
}
{ =} { \left(Dh  \right)_{v(t)}  \left(  G(v(t))  \right)
}
{ =} { \left\langle
\operatorname{Grad} \,  h  ( v(t) )  ,  G(v(t))    \right\rangle
}
}
{
\vergleichskettefortsetzungalign
{ =} { 0
}
{ } {}
{ } {}
{ } {}
}
{}{.}
Oleh karena itu, $h$ konstan pada citra $v( I)$ dan, karena

\mavergleichskette
{\vergleichskette
{ h(v(t_0))
}
{ = }{ c
}
{  }{
}
{    }{
}
{   }{
}
}
{}{}{}
berlaku

\mavergleichskette
{\vergleichskette
{ h(v(t))
}
{ = }{ c
}
{  }{
}
{    }{
}
{   }{
}
}
{}{}{}
untuk setiap

\mavergleichskette
{\vergleichskette
{ t
}
{ \in }{ I
}
{  }{
}
{    }{
}
{   }{
}
}
{}{}{,}
sehingga

\mavergleichskette
{\vergleichskette
{ v(t)
}
{ \in }{ Y
}
{  }{
}
{    }{
}
{   }{
}
}
{}{}{}
untuk setiap

\mavergleichskette
{\vergleichskette
{ t
}
{ \in }{ I
}
{  }{
}
{    }{
}
{   }{
}
}
{}{}{.}

 }

__NOEDITSECTION__

\inputaufgabeklausurloesung
{4}

{

Misalkan
\maabb {h} { I } { V
} {}
sebuah kurva yang terdiferensialkan dua kali dalam
\definitionsverweis {ruang vektor Euklides}{}{}
$V$. Buktikan bahwa jika

\mavergleichskette
{\vergleichskette
{h'(t)
}
{ \neq }{ 0
}
{  }{
}
{    }{
}
{   }{
}
}
{}{}{,}
maka berlaku persamaan

\mavergleichskettedisp
{\vergleichskette
{ \Vert { h'(t)} \Vert'
}
{ =} { { \frac{   \left\langle h'(t) , h^{\prime \prime} (t)  \right\rangle  }{  \left\langle h'(t) , h'(t)  \right\rangle  } } \Vert { h'(t)} \Vert
}
{ } {
}
{ } {
}
{ } {
}
}
{}{}{.}

}
{

Misalkan

\mavergleichskette
{\vergleichskette
{u_1
}
{ = }{ h'(t)
}
{  }{
}
{    }{
}
{   }{
}
}
{}{}{,}
yang kita lengkapi menjadi suatu basis ortogonal
\mathl{u_1,u_2 , \ldots , u_n}{} dari $V$. Misalkan
\mathl{h_1 , \ldots , h_n}{} adalah fungsi-fungsi koordinat dari $h$ terhadap basis ini. Maka

\mavergleichskettealign
{\vergleichskettealign
{ \Vert { h'(t)} \Vert'
}
{ =} { \Vert { \begin{pmatrix} h_1'(t) \\\vdots\\ h_n'(t)   \end{pmatrix} } \Vert'
}
{ =} { \sqrt{  h_1'(t)^2 + \cdots + h_n'(t)^2 }'
}
{ =} { { \frac{   h_1(t) h_1'(t) + \cdots +  h_n(t) h_n'(t) }{  \sqrt{  h_1'(t)^2 + \cdots + h_n'(t)^2 }  } }
}
{ =} { { \frac{   h_1(t) h_1'(t) + \cdots +  h_n(t) h_n'(t) }{   h_1'(t)^2 + \cdots + h_n'(t)^2  } }  \sqrt{  h_1'(t)^2 + \cdots + h_n'(t)^2 }
}
}
{
\vergleichskettefortsetzungalign
{ =} { { \frac{   \left\langle h'(t) , h^{\prime \prime} (t)  \right\rangle   }{ \left\langle h'(t) , h'(t)  \right\rangle  } } \Vert { h'(t)} \Vert
}
{ } {
}
{ } {}
{ } {}
}
{}{.}

 }

__NOEDITSECTION__

\inputaufgabeklausurloesung
{0}

{

}
{  }

__NOEDITSECTION__

\inputaufgabeklausurloesung
{0}

{

}
{  }

__NOEDITSECTION__

\inputaufgabeklausurloesung
{0}

{

}
{  }

__NOEDITSECTION__

\inputaufgabeklausurloesung
{5}

{

Buktikan pernyataan bahwa ekuivalensi tangensial lintasan-lintasan pada suatu manifold di suatu titik $P$ dapat diperiksa menggunakan sebarang peta.

}
{

Untuk suatu
\definitionsverweis {kurva diferensiabel}{}{}
\maabbdisp {\gamma} { I } { M
} {}
dengan

\mavergleichskette
{\vergleichskette
{ 0
}
{ \in }{ I
}
{  }{
}
{    }{
}
{   }{
}
}
{}{}{}
dan

\mavergleichskette
{\vergleichskette
{ \gamma(0)
}
{ = }{ P
}
{  }{
}
{    }{
}
{   }{
}
}
{}{}{}
serta suatu
\definitionsverweis {peta}{}{}
\maabbdisp {\alpha} { U } { V
} {}
\zusatzklammer {dengan

\mavergleichskettek
{\vergleichskettek
{P
}
{ \in }{ U
}
{  }{
}
{    }{
}
{   }{
}
}
{}{}{}
dan

\mavergleichskette
{\vergleichskette
{V
}
{ \subseteq }{\R^n
}
{  }{
}
{    }{
}
{   }{
}
}
{}{}{}} {} {}
ekspresi

\mathdisp {\left(  \alpha \circ \left(  \gamma \vert_{\gamma^{-1} (U)}  \right)  \right)'(0)} {  }

tidak berubah jika kita beralih ke suatu interval terbuka yang lebih kecil

\mavergleichskette
{\vergleichskette
{ 0
}
{ \in }{ I'
}
{ \subseteq }{ I
}
{    }{
}
{   }{
}
}
{}{}{}
dan suatu himpunan terbuka yang lebih kecil

\mavergleichskette
{\vergleichskette
{P
}
{ \in }{ U'
}
{ \subseteq }{ U
}
{    }{
}
{   }{
}
}
{}{}{}
\zusatzklammer {dengan peta terinduksi} {} {.}
Jadi, kita dapat mengasumsikan bahwa
\mathkor {} {\gamma_1} {dan} {\gamma_2} {}
didefinisikan pada interval yang sama dan citranya termuat dalam $U$, serta bahwa pada $U$ terdapat dua peta
\maabbdisp {\alpha_1} { U } { V_1
} {}
dan
\maabbdisp {\alpha_2} { U } { V_2
} {}
tersebut. Dari

\mavergleichskettedisp
{\vergleichskette
{ ( \alpha_1  \circ \gamma_1 )'(0)
}
{ =} { ( \alpha_1  \circ \gamma_2 )'(0)
}
{ } {
}
{ } {
}
{ } {
}
}
{}{}{}
dan
[[Totale Differenzierbarkeit/K/Kettenregel/Fakt|Kettenregel]],
dengan menggunakan sifat diferensiabel dari
\definitionsverweis {pemetaan transisi}{}{}

\mathl{\alpha_2 \circ \alpha_1^{-1}}{}, langsung diperoleh

\mavergleichskettealign
{\vergleichskettealign
{ \left(  \alpha_2 \circ \gamma_1  \right)'(0)
}
{ =} { \left(  \left(  \alpha_2 \circ \alpha_1^{-1}  \right) \circ \left(  \alpha_1 \circ \gamma_1  \right)  \right)'(0)
}
{ =} { \left(  D   \left(   \alpha_2 \circ \alpha_1^{-1}   \right)      \right)_{\alpha_1(P)} \left(   \left(  \alpha_1 \circ \gamma_1  \right)'(0)   \right)
}
{ =} { \left(  D   \left(   \alpha_2 \circ \alpha_1^{-1}   \right)      \right)_{\alpha_1(P)} \left(   \left(  \alpha_1 \circ \gamma_2  \right)'(0)   \right)
}
{ =} { \left(  \alpha_2 \circ \gamma_2  \right)'(0)
}
}
{}
{}{.}

 }

__NOEDITSECTION__

\inputaufgabeklausurloesung
{15 (3+3+4+5)}

{

Misalkan

\mavergleichskette
{\vergleichskette
{T
}
{ = }{S^1 \times S^1
}
{  }{
}
{    }{
}
{   }{
}
}
{}{}{}
merupakan torus dan

\mavergleichskette
{\vergleichskette
{S^2
}
{ \subset }{\R^3
}
{  }{
}
{    }{
}
{   }{
}
}
{}{}{}
merupakan
\definitionsverweis {permukaan bola satuan}{}{.}

a) Tunjukkan bahwa
\maabbeledisp {\varphi} {S^1 \times S^1} { S^2
} { \begin{pmatrix}  \alpha \\ \beta   \end{pmatrix} } { { \frac{   1 }{  2 } } \begin{pmatrix}  \sqrt{ 2-2 \sin \alpha   } \cos \beta  \\ \cos \alpha + \sin \beta \cos \alpha \\  1  + \sin \alpha -  \sin \beta + \sin \alpha \sin \beta   \end{pmatrix}
} {}
mendefinisikan suatu pemetaan kontinu.

b) Tunjukkan bahwa $\varphi$ surjektif.

c) Deskripsikan serat-serat $\varphi$.

d) Jelaskan pemetaan $\varphi$ dengan bantuan sebuah sketsa.

}
{

a) Pertama-tama, kita tunjukkan bahwa pemetaan tersebut benar-benar bernilai pada permukaan bola; yakni, kita harus menunjukkan bahwa jumlah kuadrat ketiga entrinya sama dengan $1$. Tanpa memperhitungkan faktor di depan, diperoleh

\mavergleichskettealign
{\vergleichskettealign
{
}
{ =} { 2  \cos^{ 2  }  \beta  -2 \sin \alpha  \cos^{ 2  }  \beta + \cos^{ 2  }  \alpha+ \cos^{ 2  }  \alpha \sin^{ 2  }  \beta +2 \cos^{ 2  }  \alpha\sin  \beta +1 + \sin^{ 2  } \alpha + \sin^{ 2  } \beta   + \sin^{ 2  } \alpha \sin^{ 2  } \beta +2 \sin \alpha  -2 \sin \beta   +2 \sin \alpha \sin \beta   -2 \sin \alpha \sin \beta +  2 \sin^{ 2  } \alpha \sin \beta  -2 \sin \alpha \sin^{ 2  } \beta
}
{ =} { 3 +   \cos^{ 2  }  \beta  -2 \sin \alpha  \cos^{ 2  }  \beta + \cos^{ 2  }  \alpha \sin^{ 2  }  \beta +2 \cos^{ 2  }  \alpha\sin  \beta  + \sin^{ 2  } \alpha \sin^{ 2  } \beta +2 \sin \alpha  -2 \sin \beta   +2 \sin^{ 2  } \alpha \sin \beta   - 2 \sin \alpha \sin^{ 2  } \beta
}
{ =} { 3 +   \cos^{ 2  }  \beta  +2 \sin \alpha   \left(  1 - \cos^{ 2  }  \beta   - \sin^{ 2  } \beta   \right)     + 2 \sin \beta \left(  - 1 + \cos^{ 2  }  \alpha    + \sin^{ 2  } \alpha  \right) + \cos^{ 2  }  \alpha\sin^{ 2  }  \beta  + \sin^{ 2  }  \alpha\sin^{ 2  }  \beta
}
{ =} { 3 +   \cos^{ 2  }  \beta  + \sin^{ 2  }  \beta
}
}
{
\vergleichskettefortsetzungalign
{ =} { 4
}
{ } {}
{ } {}
{ } {}
}
{}{.}

c) Perhitungan serat di atas kutub utara
\mathl{\begin{pmatrix}  0 \\ 0\\  1   \end{pmatrix}}{} menghasilkan syarat-syarat

\mavergleichskettedisp
{\vergleichskette
{(1 - \sin \alpha ) \cos \beta
}
{ =} { 0
}
{ =} { \cos \alpha  (1 + \sin \beta )
}
{ } {
}
{ } {
}
}
{}{}{.}
Oleh karena itu, salah satu faktor pada setiap persamaan harus sama dengan $0$. Untuk

\mavergleichskettedisp
{\vergleichskette
{ \alpha
}
{ =} { { \frac{   1 }{  2 } } \pi
}
{ } {
}
{ } {
}
{ } {
}
}
{}{}{}
kedua syarat dipenuhi untuk sebarang $\beta$. Demikian pula, untuk

\mavergleichskettedisp
{\vergleichskette
{\beta
}
{ =} { { \frac{   3 }{  2 } } \pi
}
{ } {
}
{ } {
}
{ } {
}
}
{}{}{}
kedua syarat dipenuhi untuk sebarang $\alpha$. Dalam kasus-kasus tersebut, berdasarkan

\mavergleichskettedisp
{\vergleichskette
{ \sin  { \frac{   1 }{  2 } } \pi
}
{ =} { 1
}
{ } {
}
{ } {
}
{ } {
}
}
{}{}{}
atau berdasarkan

\mavergleichskettedisp
{\vergleichskette
{ \sin  { \frac{   3 }{  2 } } \pi
}
{ =} { -1
}
{ } {
}
{ } {
}
{ } {
}
}
{}{}{}
komponen ketiga sama dengan $1$. Jika

\mavergleichskette
{\vergleichskette
{\alpha
}
{ \neq }{ { \frac{   1 }{  2 } }  \pi
}
{  }{
}
{    }{
}
{   }{
}
}
{}{}{}
dan

\mavergleichskette
{\vergleichskette
{\beta
}
{ \neq }{ { \frac{   3 }{  2 } } \pi
}
{  }{
}
{    }{
}
{   }{
}
}
{}{}{}
syarat

\mavergleichskettedisp
{\vergleichskette
{ \cos \alpha
}
{ =} { 0
}
{ =} { \cos \beta
}
{ } {
}
{ } {
}
}
{}{}{}
harus dipenuhi. Kemungkinan yang tersisa adalah

\mavergleichskette
{\vergleichskette
{ \alpha
}
{ = }{ { \frac{   3 }{  2 } } \pi
}
{  }{
}
{    }{
}
{   }{
}
}
{}{}{}
dan

\mavergleichskette
{\vergleichskette
{ \beta
}
{ = }{ { \frac{   1 }{  2 } } \pi
}
{  }{
}
{    }{
}
{   }{
}
}
{}{}{,}
tetapi dalam kasus ini komponen ketiga sama dengan $-1$.

d) Sudut $\alpha$ menentukan titik

\mavergleichskettedisp
{\vergleichskette
{P(\alpha)
}
{ =} { \begin{pmatrix}  0 \\ \cos \alpha \\  \sin \alpha   \end{pmatrix}
}
{ } {
}
{ } {
}
{ } {
}
}
{}{}{}
pada lingkaran besar yang diberikan oleh

\mavergleichskette
{\vergleichskette
{ x
}
{ = }{ 0
}
{  }{
}
{    }{
}
{   }{
}
}
{}{}{}
. Kutub utara dan
\mathl{P(\alpha)}{} menentukan titik tengah

\mavergleichskettealign
{\vergleichskettealign
{Z(\alpha)
}
{ \defeq} { \begin{pmatrix}  0 \\ 0\\  1   \end{pmatrix} + { \frac{   1 }{  2 } }  \left(  P(\alpha)   - \begin{pmatrix}  0 \\ 0\\  1   \end{pmatrix}   \right)
}
{ =} { \begin{pmatrix}  0 \\ 0\\  1   \end{pmatrix} + { \frac{   1 }{  2 } }  \begin{pmatrix}  0 \\\cos \alpha \\ \sin \alpha - 1   \end{pmatrix}
}
{ =} { { \frac{   1 }{  2 } }  \begin{pmatrix}  0 \\\cos \alpha \\  1 + \sin \alpha   \end{pmatrix}
}
{ } {
}
}
{}
{}{.}
Titik citra
\mathl{P(\alpha, \beta)}{} ditempatkan pada lingkaran $C$ dengan pusat
\mathl{Z(\alpha)}{} yang melalui kutub utara dan $P(\alpha)$ serta seluruhnya terletak pada permukaan bola. Oleh karena itu, $C$ harus terletak pada bidang yang direntang oleh

\mavergleichskettedisp
{\vergleichskette
{P(\alpha) - \begin{pmatrix}  0 \\ 0\\  1   \end{pmatrix}
}
{ =} { { \frac{   1 }{  2 } } \begin{pmatrix}  0 \\\cos \alpha \\ \sin \alpha - 1   \end{pmatrix}
}
{ } {
}
{ } {
}
{ } {
}
}
{}{}{}
dan
\mathl{\begin{pmatrix}  1 \\ 0\\  0   \end{pmatrix}}{.} Untuk menggunakan parametrisasi trigonometrik dari $C$, diperlukan suatu basis ortonormal. Oleh karena itu, kita gunakan

\mathdisp {O(\alpha) =  { \frac{   1 }{  \sqrt{ 2-2  \sin \alpha   }  } } \begin{pmatrix}  0 \\\cos \alpha \\ \sin \alpha - 1   \end{pmatrix}} { . }

Secara keseluruhan, diperoleh

\mavergleichskettealign
{\vergleichskettealign
{\varphi(\alpha)
}
{ =} { Z (\alpha) + \cos  \beta \begin{pmatrix}  1 \\ 0\\  0   \end{pmatrix}  + \sin  \beta O(\alpha)
}
{ =} {  { \frac{   1 }{  2 } }  \begin{pmatrix}  0 \\\cos \alpha \\  1 + \sin \alpha   \end{pmatrix} + \cos  \beta { \frac{   \sqrt{ 2-2  \sin \alpha     }  }{  2  } } \begin{pmatrix}  1 \\ 0\\  0   \end{pmatrix}  + \sin  \beta { \frac{   \sqrt{ 2-2  \sin \alpha     }  }{  2  } } { \frac{   1 }{  \sqrt{ 2-2  \sin \alpha   }  } } \begin{pmatrix}  0 \\\cos \alpha \\ \sin \alpha - 1   \end{pmatrix}
}
{ =} { { \frac{   1 }{  2 } }       \begin{pmatrix}        \sqrt{ 2-2  \sin \alpha     }  \cos \beta         \\\cos \alpha \sin \beta   \\    1 + \sin \alpha  - \sin \beta  + \sin \alpha \sin \beta   \end{pmatrix}
}
{ } {
}
}
{}
{}{.}

 }

__NOEDITSECTION__

\inputaufgabeklausurloesung
{5}

{

Hitunglah integral lintasan
\mathl{\int_\gamma \omega}{} untuk
\maabbeledisp {\gamma} { [-1,0] } {\R^3
} { t } {(-t^2,t^3-1,t+2)
} {,}
dengan bentuk diferensial berderajat $1$

\mavergleichskettedisp
{\vergleichskette
{ \omega
}
{ =} { x^3dx -yzdy +xz^2 dz
}
{ } {
}
{ } {
}
{ } {
}
}
{}{}{}
pada $\R^3$.

}
{

Diperoleh

\mavergleichskettealign
{\vergleichskettealign
{ \gamma^* \omega
}
{ =} { (-t^2)^3 d(-t^2) - (t^3-1) (t+2)d (t^3-1)  - t^2(t+2)^2d(t+2)
}
{ =} { 2 t^7 dt     - 3t^2 (t^4 +2t^3 -t -2) dt - t^2(  t^2  +4 t +4) dt
}
{ =} {  (2t^7 -3t^6 -6t^5   +3t^3 +6t^2 -t^4  -4t^3 -4t^2       )dt
}
{ =} { (2t^7 -3t^6 -6t^5   -t^4  -t^3 +2 t^2       )dt
}
}
{}
{}{.}
Oleh karena itu,

\mavergleichskettealign
{\vergleichskettealign
{ \int_\gamma \omega
}
{ =} { \int_{-1}^0    (2t^7 -3t^6 -6t^5   -t^4  -t^3 +2 t^2       )dt
}
{ =} { ({ \frac{   1 }{  4 } }  t^8 - { \frac{   3 }{  7 } }  t^7 - t^6- { \frac{   1 }{  5 } } t^5 - { \frac{   1 }{  4 } } t^4 + { \frac{   2 }{  3 } } t^3) \vert_{  -1 }^{  0 }
}
{ =} { -( { \frac{   1 }{  4 } }  + { \frac{   3 }{  7 } }  -1 + { \frac{   1 }{  5 } }  - { \frac{   1 }{  4 } }  - { \frac{   2 }{  3 } } )
}
{ =} { - { \frac{   3 }{  7 } }  +1 - { \frac{   1 }{  5 } } + { \frac{   2 }{  3 } }
}
}
{
\vergleichskettefortsetzungalign
{ =} { { \frac{   -45 +105-21+70 }{  105 } }
}
{ =} { { \frac{   109 }{  105 } }
}
{ } {}
{ } {}
}
{}{.}

 }

__NOEDITSECTION__

\inputaufgabeklausurloesung
{0}

{

}
{  }

__NOEDITSECTION__

\inputaufgabeklausurloesung
{3 (2+1)}

{

Misalkan
\maabbeledisp {f} {\R} {\R
} { t } {f(t)
} {,}
diberikan oleh

\mavergleichskettedisp
{\vergleichskette
{ f(t)
}
{ =} { { \frac{   \sin^{ 3  } (t^4)  }{  1+t^2 } }
}
{ } {
}
{ } {
}
{ } {
}
}
{}{}{}

a) Hitung turunan eksterior dari $f$.

b) Hitung turunan eksterior dari $fdt$.

}
{

a) Berlaku
\mathl{df=f'dt}{,} dan menurut aturan hasil bagi diperoleh

\mavergleichskettealign
{\vergleichskettealign
{f'
}
{ =} { { \frac{   (1+t^2)  3\sin^{ 2  } (t^4) \cos (t^4) 4 t^3  - 2t \sin^{ 3  } (t^4)  }{ (1+t^2)^2 } }
}
{ =} { { \frac{   12 (t^3+t^5) \sin^{ 2  } (t^4) \cos (t^4)   - 2t \sin^{ 3  } (t^4)  }{  1+2t^2+t^4 } }
}
{ } {
}
{ } {
}
}
{}
{}{.}

b) Turunan eksterior dari $fdt$ adalah
\mathl{f'dt \wedge dt =0}{.}

 }

__NOEDITSECTION__

\inputaufgabeklausurloesung
{3}

{

Berikan suatu contoh pemetaan yang terdiferensialkan secara kontinu
\maabbdisp {\varphi} { H } { H
} {}
dari suatu setengah bidang Euklides ke dirinya sendiri, dengan sifat bahwa tepat satu titik batas $H$ dipetakan ke suatu titik batas dan semua titik lainnya dipetakan ke interior setengah bidang tersebut.

}
{

Tinjau setengah bidang positif

\mavergleichskettedisp
{\vergleichskette
{ H
}
{ =} { \left\{ (x,y)  \mid   x \geq 0        \right\}
}
{ } {
}
{ } {
}
{ } {
}
}
{}{}{}
dan pemetaan
\maabbeledisp {\varphi} { H } {H
} {(x,y)} { \left(  x+y^2  , \,   y                   \right)
} {.}
Sebagai pemetaan polinomial, $\varphi$ terdiferensialkan secara kontinu. Titik batas
\mathl{(0,0)}{} dipetakan ke titik batas
\mathl{(0,0)}{.} Pada setiap titik lain dari $H$, berlaku

\mavergleichskette
{\vergleichskette
{ x
}
{ > }{  0
}
{  }{
}
{    }{
}
{   }{
}
}
{}{}{}
atau

\mavergleichskette
{\vergleichskette
{ y
}
{ \neq }{ 0
}
{  }{
}
{    }{
}
{   }{
}
}
{}{}{}
dan oleh karena itu selalu

\mavergleichskettedisp
{\vergleichskette
{ x+y^2
}
{ >} {  0
}
{ } {
}
{ } {
}
{ } {
}
}
{}{}{,}
sehingga komponen pertamanya positif.

 }

__NOEDITSECTION__

\inputaufgabeklausurloesung
{0}

{

}
{  }

__NOEDITSECTION__

\inputaufgabeklausurloesung
{10}

{

Buktikan versi balok dari teorema Stokes.

}
{

Karena kedua ruas persamaan ini linear terhadap $\omega$, kita dapat mengasumsikan bahwa $\omega$ berbentuk

\mavergleichskettedisp
{\vergleichskette
{ \omega
}
{ =} { f dx_2 \wedge \ldots \wedge dx_n
}
{ } {
}
{ } {
}
{ } {
}
}
{}{}{}
dengan suatu lingkungan terbuka dari $Q$ tempat fungsi $f$ yang terdiferensialkan secara kontinu didefinisikan. Integral pada kedua ruas adalah
\definitionsverweis {integral Lebesgue}{}{}
dari fungsi kontinu pada himpunan bagian
\mathkor {} {\R^n} {atau} {\R^{n-1}} {.}
Oleh karena itu, pada kedua ruas kita dapat beralih ke
\definitionsverweis {penutupan topologis}{}{,}
karena hal ini hanya mengubah domain integrasi terkait sebesar suatu himpunan nol dan karenanya tidak mengubah nilai integral.

Tuliskan balok tertutup sebagai

\mavergleichskettedisp
{\vergleichskette
{ \overline{Q}
}
{ =} { [a,b] \times \tilde{Q}
}
{ } {
}
{ } {
}
{ } {
}
}
{}{}{.}
Terapkan
[[Differentialform/Lokal/Zurückziehen unter partiell konstanter Abbildung/Fakt|Korollar 14.10 (Differentialgeometrie (Osnabrück 2023))]]
pada setiap sisi $S$ kecuali
\mathkor {} {a \times \tilde{Q}} {dan} {b \times \tilde{Q}} {,}
sehingga pada sisi-sisi tersebut diperoleh

\mavergleichskettedisp
{\vergleichskette
{
\int_{ S  }   \omega
}
{ =} {
\int_{  S  }  f dx_2 \wedge \ldots \wedge dx_n
}
{ =} {
\int_{ S  }   0
}
{ =} { 0
}
{ } {
}
}
{}{}{,}
karena pada setiap sisi tersebut salah satu variabel
\mathl{x_2 , \ldots , x_n}{} konstan. Berdasarkan
[[Sigmaendliche Räume/Satz von Fubini/Fakt|Satzes von Fubini]]
dan
[[Riemann-Integral/Hauptsatz/Newton-Leibniz/Fakt|Hauptsatzes der Infinitesimalrechnung]]
\zusatzklammer {diterapkan untuk setiap \mathlk{(x_2 , \ldots , x_n)}{} yang tetap} {} {,}
berlaku

\mavergleichskettealignhandlinks
{\vergleichskettealignhandlinks
{
\int_{ Q  }   d\omega
}
{ =} {
\int_{  Q  }   df \wedge  dx_2 \wedge \ldots \wedge dx_n
}
{ =} {
\int_{  Q  }   \left(  \sum_{j = 1}^n { \frac{   \partial   f   }{  \partial   x_j   } } dx_j  \right) \wedge  dx_2 \wedge \ldots \wedge dx_n
}
{ =} {
\int_{  Q  }   { \frac{   \partial   f   }{  \partial   x_1   } } dx_1 \wedge  dx_2 \wedge \ldots \wedge dx_n
}
{ =} {
\int_{ \tilde{Q}   }   \left(  \int_{[a,b] } { \frac{   \partial   f   }{  \partial   x_1   } } dx_1  \right)  dx_2 \wedge \ldots \wedge dx_n
}
}
{
\vergleichskettefortsetzungalign
{ =} {
\int_{  \tilde{Q}   }   \left(  f(b,x_2 , \ldots , x_n ) - f(a,x_2 , \ldots , x_n )   \right) dx_2 \wedge \ldots \wedge dx_n
}
{ =} {
\int_{ b \times \tilde{Q}   }   f dx_2 \wedge \ldots \wedge dx_n  -
\int_{  a \times \tilde{Q}   }  f dx_2 \wedge \ldots \wedge dx_n
}
{ =} { \sum_{S \text{ orientierte Seite von } Q}
\int_{  S  }  f dx_2 \wedge \ldots \wedge dx_n
}
{ =} {
\int_{  \partial Q  }   f dx_2 \wedge \ldots \wedge dx_n
}
}
{
\vergleichskettefortsetzungalign
{ =} {
\int_{  \partial Q  }   \omega
}
{ } {}
{ } {}
{ } {}
}{.}

 }

__NOEDITSECTION__

\inputaufgabeklausurloesung
{4}

{

Untuk kurva
\maabbeledisp {\psi} {\R} { {\mathbb H}
} { t } { \left(  t   , \,  e^t                   \right)
} {,}
yang bernilai dalam
\definitionsverweis {setengah bidang Poincaré}{}{}
${\mathbb H}$, tentukan
\definitionsverweis {simbol-simbol Christoffel}{}{}
untuk
\definitionsverweis {koneksi tarikan balik}{}{}
\zusatzklammer {dari
\definitionsverweis {koneksi Levi-Civita}{}{}
pada ${\mathbb H}$} {} {}
di atas $\R$.

}
{

Satu-satunya turunan standar pada $\R$ adalah $\partial$, sehingga indeks bawah pada simbol-simbol Christoffel kita abaikan. Menurut
[[Linearer Zusammenhang/Vertikale Ableitung/Längs Abbildung/Eigenschaften/Fakt|Fakt]] [[Kurs:Differentialgeometrie/Linearer Zusammenhang/Vertikale Ableitung/Längs Abbildung/Eigenschaften/Fakt/Faktreferenznummer|*****]]  (2)
untuk

\mavergleichskette
{\vergleichskette
{ \ell,k
}
{ = }{ 1,2
}
{  }{
}
{    }{
}
{   }{
}
}
{}{}{}

\mavergleichskettedisp
{\vergleichskette
{ \tilde{\Gamma}_\ell^k (t)
}
{ =} { \psi_1'(t) \Gamma_{1 \ell}^k (\psi(t)) + \psi'_2(t) \Gamma_{2 \ell}^k (\psi(t))
}
{ =} { \Gamma_{1 \ell}^k (t,e^t)   + e^t \Gamma_{2 \ell}^k (t,e^t)
}
{ } {}
{ } {}
}
{}{}{.}
Dengan
[[Halbebene/Poincaré/Riemannsche Geometrie/Christoffelsymbole/Beispiel|Beispiel]] [[Kurs:Differentialgeometrie/Halbebene/Poincaré/Riemannsche Geometrie/Christoffelsymbole/Beispiel/Beispielreferenznummer|*****]]
diperoleh

\mavergleichskettedisp
{\vergleichskette
{ \tilde{\Gamma}_1^1 (t)
}
{ =} { \Gamma_{1 1}^1 (t,e^t) + e^t \Gamma_{2 1}^1 (t,e^t)
}
{ =} {  -e^t e^{-t}
}
{ =} { -1
}
{ } {
}
}
{}{}{,}

\mavergleichskettedisp
{\vergleichskette
{ \tilde{\Gamma}_1^2 (t)
}
{ =} { \Gamma_{1 1}^2 (t,e^t) + e^t \Gamma_{2 1}^2 (t,e^t)
}
{ =} { e^{-t}
}
{ } {
}
{ } {
}
}
{}{}{,}

\mavergleichskettedisp
{\vergleichskette
{ \tilde{\Gamma}_2^1 (t)
}
{ =} { \Gamma_{1 2}^1 (t,e^t) + e^t \Gamma_{2 2}^1 (t,e^t)
}
{ =} {  -e^{-t}
}
{ } {
}
{ } {
}
}
{}{}{}
dan

\mavergleichskettedisp
{\vergleichskette
{ \tilde{\Gamma}_2^2 (t)
}
{ =} { \Gamma_{1 2}^2 (t,e^t) + e^t \Gamma_{2 2}^2 (t,e^t)
}
{ =} { - e^t e^{-t}
}
{ =} { -1
}
{ } {
}
}
{}{}{.}

 }

 [[Kategorie:Latexseite]]
